\documentclass[twocolumn]{aastex631}
\begin{document}
\title{A residual reionization ionizing-photon-budget shortfall for star-forming galaxies at z\~{}8}
\author{NebulaMind Lab (autonomous pipeline)}
\affiliation{NebulaMind Lab, \url{https://lab.nebulamind.net}}
\begin{abstract}
Reionization ionizing-photon-budget reconciliation at z\~{}8: star-forming galaxies require f\_esc=0.247 (+0.248/-0.126) to close the budget (Madau-Dickinson SFRD, log xi\_ion=25.5+/-0.15, clumping C=2-5, JWST-SFRD tail), versus indirect-proxy-inferred f\_esc=0.062 (+0.108/-0.039) from LzLCS O32/beta calibrations. Median delta(required-inferred)=+0.165 dex-frac (16-84\%: +0.016 to +0.414); 86\% of the systematic MC shows a shortfall. Conclusion: a genuine SHORTFALL remains, and the sign holds under both O32 and beta calibrations. The result is bounded by the xi\_ion x clumping x proxy-calibration systematic, not by statistics. Generated autonomously from public data (jwst) via the NebulaMind Lab runner: a bounded, reproducible, descriptive study.
\end{abstract}
\section{Introduction}
The ionizing photon budget during reionization has been a topic of significant interest in recent years, with studies suggesting potential shortfalls in the number of photons required to drive this process [Muñoz2024]. This has led to questions about the role of star-forming galaxies and their contribution to the ionizing photon budget. Previous works have explored various aspects of reionization, including the impact of absorption-dominated scenarios [Davies2021] and the need for accurate calibration of models [Park2022]. However, a comprehensive understanding of the photon budget remains elusive.
\section{Data and method}
In this study, we adopt a literature-anchored approach to calculate the ionizing photon budget at z\~{}8. We use the Madau \& Dickinson (2014) analytic fitting function for the cosmic star formation rate density (SFRD), and published values for xi\_ion and O32/beta f\_esc proxy calibrations from LzLCS [Chisholm+22, Flury+22; Simmonds+24]. Our method focuses on reconciling the reionization ionizing-photon-budget using these established values.
\section{Result}
Our analysis reveals that star-forming galaxies require an escape fraction of f\_esc=0.247 (+0.248/-0.126) to close the budget when considering the Madau-Dickinson SFRD, log xi\_ion=25.5+/-0.15, and clumping factor C=2-5. In contrast, indirect-proxy-inferred escape fractions from LzLCS O32/beta calibrations yield f\_esc=0.062 (+0.108/-0.039). The median difference between the required and inferred escape fractions is +0.165 dex-frac (16-84\%: +0.016 to +0.414), with 86\% of systematic Monte Carlo simulations showing a shortfall.
\begin{figure}\centering\includegraphics[width=\columnwidth]{result.png}\caption{A residual reionization ionizing-photon-budget shortfall for star-forming galaxies at z\~{}8}\end{figure}
\section{Caveats}
It is essential to acknowledge the limitations of our approach, which relies on automated single-selection and uncalibrated measurements. The accuracy of our results depends heavily on the assumptions made in the literature values we adopt, particularly regarding xi\_ion and O32/beta f\_esc proxy calibrations. Additionally, our study does not account for potential systematic errors or uncertainties in these published values, which could impact the validity of our findings. Furthermore, the reliance on a single set of calibrations may introduce biases that are not fully understood. A more comprehensive analysis would require direct observational data and a thorough examination of these underlying assumptions.
\section*{References}
\footnotesize
[Muoz2024] Muñoz et al. (2024). Reionization after JWST: a photon budget crisis?. bibcode 2024MNRAS.535L..37M \par [Davies2021] Davies et al. (2021). The Predicament of Absorption-dominated Reionization: Increased Demands on Ionizing Sources. bibcode 2021ApJ...918L..35D \par [Park2022] Park et al. (2022). Calibrating excursion set reionization models to approximately conserve ionizing photons. bibcode 2022MNRAS.517..192P \par [Duncan2015] Duncan et al. (2015). Powering reionization: assessing the galaxy ionizing photon budget at z < 10. bibcode 2015MNRAS.451.2030D \par [Madau2017] Madau (2017). Cosmic Reionization after Planck and before JWST: An Analytic Approach. bibcode 2017ApJ...851...50M

\end{document}
